\documentclass[11pt,a4paper]{article}
\usepackage[T1]{fontenc}
\usepackage{amsmath,amssymb,booktabs,longtable,geometry,hyperref}
\geometry{margin=26mm}
\title{EPPA--III: suggested manuscript corrections}
\author{Independent validation register}
\date{10 October 2026}
\begin{document}
\maketitle

\noindent\textbf{Source and policy.}
This is a register of proposed corrections to the manuscript
\emph{Graphs with Small EPPA Witnesses}, referenced against the
user-supplied snapshot \texttt{main-new(7).tex}.
All anchors below use labels present in that source.
The original manuscript is \textbf{frozen}: neither this report nor
passing GitHub Actions authorizes an edit to its TeX.

\paragraph{Evidence levels.}
\textbf{Confirmed}: mathematical counterexample, exact definition
comparison, or unambiguous textual error.
\textbf{Paper-level}: a proposed mathematical proof, not yet
completely certified in Lean.
\textbf{Pending}: a proof dependency or author decision remains.
The Lean CI build certifies only declarations explicitly included
in \texttt{CheckAxioms.lean}; small exhaustive checks do not
establish arbitrary-size results.

\section*{Issue index}
\begin{longtable}{@{}p{15mm}p{66mm}p{32mm}p{35mm}@{}}
\toprule
ID & Anchor in supplied source & Evidence & Proposed action\\
\midrule
C01 & After \texttt{def:generalizedcover} & Confirmed & Taylor-double example\\
C02 & \texttt{claim:taylorconclusion} & Confirmed & Mate-edge qualification\\
C03 & \texttt{thm:sKt_witnesses}, alternative (2) & Confirmed & Wagner complement\\
C04 & \texttt{claim:connections} & Paper-level & Orbit counting\\
C05 & Claim 3.11 in \texttt{thm:sKt_witnesses} & Paper + finite & Wagner proof\\
C06 & Immediately before Claim 3.11 & Confirmed & Missing brace\\
C07 & \texttt{lemma:different_cliques} & Confirmed & Hypothesis\\
C08 & \texttt{prop:imprimitive_summary} & Confirmed & Restore \(2G\)\\
C09 & End of \texttt{prop:imprimitive_summary} & Confirmed gap & Complement case\\
C10 & Use of \texttt{thm:sKt_witnesses} & Paper-level & Partition lemma\\
C11 & \texttt{lem:C5_10vertex} & Paper-level & Cross cycles\\
C12 & \texttt{lem:R3_18vertex} & Paper + finite & Orbit proof\\
C13 & Claims 3.8, 3.10, 3.15, 3.16 & Pending & Proof alternatives\\
\bottomrule
\end{longtable}

\section*{Confirmed localized corrections}

\subsection*{C01. Taylor double of a complete graph}
\textbf{Location:} approximately source lines 522--523, immediately
after Definition~\texttt{def:generalizedcover}.
Definition~\texttt{def:doubles} makes the Taylor double of
\(K_t\) equal to \(2K_t\), with no cross-edges. The connection
type \(M\) adds a matching, and \(\overline M\) is not the
graph complement of that Taylor double.

\textbf{Proposed two-sentence replacement:}
\begin{quote}
For \(G\cong K_t\), the Taylor double is \(2K_t\),
corresponding to connection type \(X=0\).
Connection type \(X=M\) instead gives the Taylor double
with an additional matching between corresponding vertices;
this is another generalised clique cover.
\end{quote}
Keep the next paragraph about transversals of the matching
edges as referring to \(X=M\).

\subsection*{C02. Claim 3.6 and the mate matching}
\textbf{Location:} \texttt{claim:taylorconclusion}.
The argument concerns edges between distinct fibres, not
the edge inside each fibre. For
\(G=C_5\dot\cup K_1\), adding all mate edges to the
Taylor double yields a genuine transitive imprimitive EPPA
witness for \(G\) (see the exact finite test
\texttt{checks/section3\_mate\_matching.py}).

\textbf{Replacement statement:}
\begin{quote}
The edges of \(H\) between distinct blocks agree with those
of the Taylor double of \(G\). Within each block, a
uniform matching may additionally be present.
\end{quote}
This makes explicit the existing qualification
\emph{up to complementation within all blocks} in
Lemma~\texttt{lemma:doublecover}.

\subsection*{C03. Wagner complement in Lemma 3.7}
\textbf{Location:} alternative (2) of
\texttt{thm:sKt_witnesses}, approximately source line 1225.

\textbf{Action:} delete \emph{or its complement} after
\emph{the Wagner graph}. The displayed matrix gives a
3-regular graph; its complement is 4-regular.
The complemented classification is obtained by
complementing the pair \((G,H)\), not by relabelling
the same displayed matrix.

\subsection*{C06. Missing brace before Claim 3.11}
\textbf{Location:} approximately source line 1310.

\textbf{Action:} replace the malformed
\(Y\in\{M,\overline M\) with
\(Y\in\{M,\overline M\}\).
No mathematical change is involved.

\end{document}
